\section*{\texorpdfstring{\S\CurrentExerciseGroup}{第{\CurrentExerciseGroup}节}}
\phantomsection
\addcontentsline{toc}{section}{第{\CurrentExerciseGroup}节习题}

\begin{enumerate}
\def\labelenumi{\arabic{enumi})}
\item
  设 E 是一个具有可数标准正交基 \((\mathbf{e}_{n})\) 的 Hilbert 空间（TVS, V, §2, No.3），\(E_{0}\) 是 E 中由有限个向量 \(e_{n}\) 的线性组合组成的稠密线性子空间。设 X 是 R 中由 0 及诸点 1/n （n 为整数 \(\geqslant 1\) ）组成的紧子空间；设 \(\mu\) 是 X 上由每个点 1/n 处的质量 \(1/n^{2}\) 及点 0 处的质量 0 定义的正测度（§1, No.~3, 例 I）。考虑 X 到 \(E_{0}\) 中的连续映射 f ，它由 \(\mathbf{f}(0)=0\) ，\(\mathbf{f}(1/n)=\mathbf{e}_{n}/n\) （对 \(n\geqslant1\) ）定义；试证积分 \(\int f d\mu\) 不属于 \(E_{0}\) 。
\item
  设空间 X 、E 与映射 f 如习题 1 中所定义，试证：映射 \(\lambda \mapsto \lambda(\mathbf{f})\) （从 \(\mathcal{M}(\mathrm{X};\mathbf{C})\) 到 E 中）关于淡拓扑不连续。（若 \(g_{k}\) （ \(1 \leqslant k \leqslant n\) ）是 X 上的 \(n\) 个连续标量函数，试证：存在 X 上的离散测度 \(\lambda\) ，使 \(\int g_{k} d\lambda = 0\) 对 \(1 \leqslant k \leqslant n\) 成立，但 \(\| \int \mathbf{f} d\lambda \|\) 可任意大。）
\item
  \begin{enumerate}
  \def\labelenumii{\alph{enumii})}
  \tightlist
  \item
    设 E 是拟完备的 Hausdorff 局部凸空间。试证：当 \(\mathcal{M}(\mathrm{X};\mathbf{C})\) 赋以拟强拓扑（§1，习题 8）（相应地，严格紧收敛拓扑）时，映射 \((\mathbf{f},\mu)\mapsto\int\mathbf{f}d\mu\) 是 \((\mathfrak{S},\mathfrak{T})\) -亚连续的，其中 \(\mathfrak{T}\) 是 \(\mathcal{M}(\mathrm{X};\mathbf{C})\) 的淡有界子集所成之集，\(\mathfrak{S}\) 是 \(\mathcal{K}(\mathrm{X};\mathrm{E})\) 中含于 \(\mathcal{K}(\mathrm{X},\mathrm{K};\mathrm{E})\) （K 为适当的紧集）的有界（相应地，紧）子集所成之集。
  \end{enumerate}
\end{enumerate}

\begin{enumerate}
\def\labelenumi{\alph{enumi})}
\setcounter{enumi}{1}
\tightlist
\item
  对每个紧子集 \(\mathbf{K}\) ⊂ \(\mathbf{X}\) ，映射 \((\mathbf{f},\mu)\mapsto \int \mathbf{f}d\mu\) 在 \(\mathcal{H}(\mathrm{X},\mathrm{K};\mathrm{E})\times \mathcal{M}(\mathrm{X};\mathbf{C})\) 上连续（当 \(\mathcal{M}(\mathrm{X};\mathbf{C})\) 赋以拟强拓扑时），且在 \(\mathcal{H}(\mathrm{X},\mathrm{K};\mathrm{E})\times \mathcal{M}_{+}(\mathrm{X})\) 上连续（当 \(\mathcal{M}_{+}(\mathrm{X})\) 赋以淡拓扑时）。
\end{enumerate}

\begin{enumerate}
\def\labelenumi{\arabic{enumi})}
\setcounter{enumi}{3}
\tightlist
\item
  利用 No.~4 的命题 8 证明 No.~2 的命题 4 的推论（化归到 E 拟完备的情形，并利用 §1, No.~9, 命题 15 的推论 3）。
\end{enumerate}

\setcounter{exercisegroup}{4}
\edef\CurrentExerciseGroup{\arabic{exercisegroup}}
